% OWNER: sections/s4_4_validity.tex. Edit only this file.
\seclabel
\claim{\secnum{4.4}{sec:validity}\quad Tests on one account on one morning limit these results; five questions remain open}

\lead{Every measurement in this report comes from one account on one morning, 05:38 to 10:34 KST on 24 Nov 2026, apart from the browser recording of 28 Oct 2026 that gave the 3\,s poll.} All requests used one portfolio and were sent from one PC on one network, with the paths, headers and bodies the web client sends, by curl and by the wrapper's httpx client. A value that depends on load, on the day or on the network can therefore differ on another day: the fill pace of Section~\ref{sec:orders}, the reply sizes of Section~\ref{sec:feed} and every timing. We did not test request forms that the web client never sends, nor a second account.

\lead{The evidence is thinnest where a failure costs most.} The 15\,s window in which the wrapper looks for a placed order in the feed rests on the wrapper's live buys, found 1 to 3\,s after \code{AddOrder}, and on the 08:30 order, in the feed by 08:33; no sell was timed and no slow server was tested. A window that is too short turns a placed order into \code{AmbiguousState}. Connection timing, where a wrong value costs only latency, rests on 50 requests (10 runs of 5, Section~\ref{sec:transport}).

\lead{Most layer timings were taken before the 09:00 open, so market-hours load is measured only for the Python client.} The 16.0\ms{} and 106.5\ms{} of Section~\ref{sec:transport} come from curl before the open; the wrapper's 31 to 39\ms{} on a reused connection come from the session. The two differ in client, time of day and load at once, and the report does not separate these causes.

\lead{The idle timeout, the rollover, the token tolerance, \code{Logout}, the order's appearance in the feed and the resend after a 401 rest on 2 observations or fewer.} The idle timeout of 15 to 30\,s rests on 1 connection reused after 5\,s and 15\,s of silence and 1 replaced after 30\,s. The business-day rollover between 06:02:10 and 07:06:44 rests on 2 readings of \code{currBizDay} on 1 day. The 300\,s tolerance after token expiry rests on 2 tests (Section~\ref{sec:auth}), and the finding that a saved cookie jar keeps renewing after \code{Logout} rests on 1 call with 1 saved jar. The time from \code{AddOrder} to the order's appearance in the feed was 1 to 3\,s for the wrapper's live buy and is otherwise known only as an upper bound (the 08:30:05 order was in the feed by 08:33), so the 15\,s search window of Section~\ref{sec:orders} has not been checked against a slow server. The resend of an order after a 401 rests on 0 observations on an order request and on the framework's documentation alone (Section~\ref{sec:auth}). Where a wrapper parameter depends on one of them, it is set with a margin inside the observed value (Appendix~\ref{app:decisions}); none has been measured on a second day.

\lead{The server's behaviour is undocumented, so any of these findings can change without notice.} The delta rule, the meaning of \code{wei}, the rounding of \code{qty2ord} and the 300\,s tolerance were inferred from replies; the organiser has not confirmed them. The wrapper checks the ones it can at run time: it refuses orders while \code{currBizDay} differs from the KST date, raises \code{StaleFeed} when the last feed reply is older than 26\,s, and raises \code{AmbiguousState} when a placed order cannot be found. The live checks in \code{tests/live} repeat the measurements; 3 of them (\code{live_orders}, \code{live_sell} and \code{live_rounding}) place and stop real orders.

\lead{The organiser's limit applies per IP address, so club members behind one shared NAT address share it.} The organiser's reply [5] states that an IP address that places 100 or more orders in a day may be blocked, and so may bursts of traffic that put a heavy load on the system. The wrapper's cap of 80 orders a day counts 1 account only. Members who run the wrapper from the same campus network leave through one public address, so their orders and requests may count together against a limit the wrapper cannot see. The load threshold itself is unknown: the timing study sent runs of 5 sequential requests, and no load test was run.

\lead{Two guarantees hold only under an operating assumption or at one moment.} First, the order reply carries no ID, so the wrapper cannot tell its own order from an identical one placed at the same moment from the web page or another process. It raises \code{AmbiguousState} when more than one new record matches, but a single external match cannot be detected, so the wrapper assumes one submitter per account while it runs. Second, \code{plan()} checks target weights against the stock, sector and total limits with the data at planning time. Sells go out before buys but need not fill first (7 shares filled over 7\,min 37\,s, Section~\ref{sec:orders}), so a sector or the total can pass its limit while a rebalance fills; when the organiser measures violations is question Q4. \figref{guarantees} lists what a user can rely on, the assumption each guarantee needs, its evidence and what was not tested.

\lead{No fault was injected live, and the code changed after the live tests has run only offline.} No feed disconnect, reconnection or backoff and no connection cut during an \code{AddOrder} was produced against the live server, and no process has yet run past 720\,s to renew inside one run. After the live tests, 6 behaviours were changed: one reply deadline per feed request, the stored cookie expiry, the search after a failed resend, the exclusion of order IDs already returned, \code{AmbiguousState} when more than one new record matches an order, and the halt of writes after a 403 or 429 on any request. Seven of the 31 offline tests cover these changes; none has run against the live server.

\lead{Whether the stored cookie jar still renews after a password change is untested (Q3 in \figref{openq}).} Until it is, the wrapper treats \code{session.json} as a 7-day credential that keeps renewing after \code{Logout} (Section~\ref{sec:auth}); the wrapper keeps it with mode 0600 outside the repository (\figref{security}).

\lead{Our permission to automate orders and to reproduce the web client's requests rests on the organiser's e-mail of 24 Nov 2026 alone [5].} The reply states that the organiser provides no official API and does not prohibit automating orders with scripts or web drivers (Section~\ref{sec:intro}); we kept a copy. This position exists only in a reply to our enquiry, and the organiser can change it at any time.

\lead{Five questions remain open: 4 can be settled by tests and 1 only by the organiser (\figref{openq}).} For each, the wrapper follows an interim behaviour that is safe if the answer goes the other way, such as counting a stop as an order until the organiser says whether cancels count toward the 100 (Q2). Three earlier questions were settled on 24 Nov. \code{qty2ord} is rounded down: a Samsung Electronics order worth 1.76 shares asked for 1 share, twice (Section~\ref{sec:oneshare}). The state \code{Accepted} was observed right after placement (Section~\ref{sec:orders}), so Q5 now asks about \code{Generated} and \code{Deleted} and about how a windowed order that reaches \code{hm1} unfilled ends. \code{bs_updated} carries the balance fields \code{prft}, \code{na}, \code{rt}, \code{nav} and \code{stkAmt} in every delta (Section~\ref{sec:feed}).

% OWNER: figs/openq.tex. Edit only this file.
\fig[wide]{openq}{5 questions remain open: 4 can be settled by tests and 1 only by the organiser}[Open questions, what the wrapper does until each is settled, and the next step]{%
{\tabfont\begin{tabularx}{\linewidth}{>{\RaggedRight\arraybackslash}p{6mm} >{\RaggedRight\arraybackslash}X >{\RaggedRight\arraybackslash}p{50mm} >{\RaggedRight\arraybackslash}p{50mm}}
\thd{} & \thd{Question} & \thd{Interim behaviour} & \thd{How it will be settled (next step)}\\
Q1 & Do the \code{CancelOrder}, \code{ModifyOrder} and \code{DeleteOrder} bodies taken from the client code work? None has been tested live & The wrapper does not use them; it stops orders with \code{CancelOrderBatch} & Capture the requests when the web page offers them for an order\\\rowsep
Q2 & Does the 100-order limit count stops, and does it apply per IP address or per account? & Each stopped order counts as 1 order; the wrapper caps the day at 80 orders & Ask the organiser\\\rowsep
Q3 & Does a password change revoke a saved cookie jar? & The cookie jar is treated as a 7-day credential (\ref{sec:design}) & Change the password, then renew with a cookie jar saved before the change\\\rowsep
Q4 & When does the server count a sector or small-cap violation: at any moment in the trading session, or at the close? Sector limits move with market weights (\ref{sec:guidelines}) & \code{plan()} rejects sector breaches before any order exists, with limits read from the feed at each call; \code{guidelines()} reports the small-cap total & Log \code{Violations} and the feed's sector weights every hour and compare\\\rowsep
Q5 & When do the order states \code{Generated} and \code{Deleted} appear? & Every state other than \code{Done} counts as open: \code{plan()} subtracts its remaining size and \code{wait_done} keeps waiting & Log every state change in \code{ords}, including orders the server refuses\\
\end{tabularx}}}{Settled on 24 Nov and removed from this list: whether the live feed pushes (\ref{sec:feed}), the format of \code{updated} and the exactness of the delta merge (\ref{sec:feed}), parallel renewal (\ref{sec:auth}), the \code{Orders} parameters (\ref{sec:queries}), the final state of filled and stopped orders and the state \code{Accepted} right after placement (\ref{sec:orders}), the sell order form (\ref{sec:wrappertests}), the rounding of \code{qty2ord} (floor, \ref{sec:oneshare}), the content of \code{bs_updated} (\ref{sec:feed}), the source of net assets (\ref{sec:oneshare}) and the feed fields behind the guidelines (\ref{sec:guidelines}).}{DART Quant Team}

% OWNER: figs/guarantees.tex. Edit only this file.
\fig[wide]{guarantees}{What a user can rely on, and under which assumption}[Each guarantee of \code{pytimefolio} 0.1.0 with the assumption it needs, the evidence behind it and what was not tested]{%
{\tabfont\begin{tabularx}{\linewidth}{>{\RaggedRight\arraybackslash}p{38mm} >{\RaggedRight\arraybackslash}p{42mm} >{\RaggedRight\arraybackslash}X >{\RaggedRight\arraybackslash}p{44mm}}
\thd{Guarantee the user may rely on} & \thd{Assumption it needs} & \thd{Evidence} & \thd{What was not tested}\\
The merged state equals the server's snapshot with the same \code{t} & One feed request in flight; the server keeps the delta rule (a) & Live: \code{pos} and \code{ords} equal in 15 of 15 rounds, 09:35 (\ref{sec:feed}) & \code{bs} and \code{msecs} field by field; rules (b) and (c) during the session\\\rowsep
No read, order or \code{wait()} uses a state older than 26\,s & The feed task is running; a poll every 3\,s is a target, not a guarantee & Offline: a stale feed raises \code{StaleFeed}; pings do not extend the 20\,s reply deadline (\ref{sec:wrappertests}) & A live feed outage\\\rowsep
\hl An order the server may have received is never resent & A 401 comes before the order handler [9]; a connection error before set-up means the request never left & Offline: a lost reply not resent; a failed resend searched, no third send (\ref{sec:orders}) & A 401 on an order; a connection cut during a live \code{AddOrder}\\\rowsep
\hl A placed order is confirmed within 15\,s, or \code{AmbiguousState} is raised & The server shows a new order in the feed or in \code{Orders} within 15\,s & Live: buys found 1 to 3\,s after \code{AddOrder}; the 08:30 order by 08:33 (\ref{sec:validity}) & A slow server; the delay for sells\\\rowsep
\hl A returned order ID belongs to the user's order & One submitter per account while the wrapper runs & Offline: 2 identical orders get 2 IDs; 2 new matching records raise \code{AmbiguousState} (\ref{sec:orders}) & A single identical order from the web page or another process cannot be detected\\\rowsep
At most 80 orders a day, at least 1\,s between writes & One process per count file; a cancel counted as an order & Offline: cap, sync and 1.002\,s spacing; live: counter 16 of 80 by 10:34 (\ref{sec:design}) & Several processes on one count file; whether the organiser counts cancels (Q2)\\\rowsep
\hl Target weights meet the stock, sector and total limits at planning time & Feed values (\code{weiLimit}, \code{msecs}) current when \code{plan()} runs & Live: +6\,pp Industrials rejected, +3\,pp accepted; targets above 100\% rejected (\ref{sec:guidelines}) & Weights during execution: a sell sent first need not fill first; when violations are counted (Q4)\\\rowsep
All writes stop after any 403 or 429 & The server signals a block with 403 or 429 & Offline: a 403 or 429 on a write or a read halts writes; a blocked feed stops reconnecting (\ref{sec:design}) & No block was seen live\\\rowsep
Renewal without the password for up to 7 days & The server keeps the 7-day cookie expiry; renewal at 720\,s & Live: a second process renewed from the stored jar; offline: the expiry survives a restart (\ref{sec:auth}) & A process running past 720\,s or for 7 days\\
\end{tabularx}}}{Highlighted rows rest on an operating assumption or on evidence thinner than the cost of a failure. Section numbers in brackets give the evidence. \figref{goalstatus} judges the goals against the same tests.}{DART Quant Team; \figref{tests}}

